\documentclass[11pt]{article}

\usepackage{fullpage}
\usepackage{amsmath}
\usepackage{amssymb}
\usepackage{mathpartir}

\title{Lab 14: Evaluation Order and Typing Rules}
\author{Kushagra Bainsla}
\date{}

% Language forms
\newcommand{\true}{\mbox{\tt true}}
\newcommand{\false}{\mbox{\tt false}}
\newcommand{\ife}[3]{\mbox{\tt if}~{#1}~\mbox{\tt then}~{#2}~\mbox{\tt else}~{#3}}
\newcommand{\suc}[1]{\mbox{\tt succ}~{#1}}
\newcommand{\prd}[1]{\mbox{\tt pred}~{#1}}
\newcommand{\iszero}[1]{\mbox{\tt iszero}~{#1}}
\newcommand{\concat}[2]{\mbox{\tt concat}~{#1}~{#2}}
\newcommand{\isemptystr}[1]{\mbox{\tt isemptystr}~{#1}}
\newcommand{\strlen}[1]{\mbox{\tt strlen}~{#1}}

% Judgments
\newcommand{\step}[2]{{#1} \rightarrow {#2}}
\newcommand{\typeof}[2]{\vdash {#1} : {#2}}

\begin{document}
\maketitle

\section*{Language}

\[
\begin{array}{rcl}
e &::=& \true \mid \false \mid i \mid s \\
  &\mid& \suc e \mid \prd e \mid \iszero e \\
  &\mid& \ife e e e \\
  &\mid& \concat e e \mid \isemptystr e \mid \strlen e
\end{array}
\]

\[
\begin{array}{rcl}
v &::=& \true \mid \false \mid i \mid s
\end{array}
\]

\[
\begin{array}{rcl}
T &::=& \mbox{Bool} \mid \mbox{Int} \mid \mbox{String}
\end{array}
\]

\noindent
Here, $i \in \mathbb{Z}$ and $s$ is a string literal.

\section*{Small-Step Evaluation Order Rules}

\noindent
Evaluation is left-to-right and call-by-value.

\subsection*{Conditionals}
\begin{mathpar}
\inferrule
  {\step{e_1}{e_1'}}
  {\step{\ife{e_1}{e_2}{e_3}}{\ife{e_1'}{e_2}{e_3}}}
\quad\textsc{(E-If)}

\inferrule
  { }
  {\step{\ife{\true}{e_2}{e_3}}{e_2}}
\quad\textsc{(E-IfTrue)}

\inferrule
  { }
  {\step{\ife{\false}{e_2}{e_3}}{e_3}}
\quad\textsc{(E-IfFalse)}
\end{mathpar}

\subsection*{Integer Operators}
\begin{mathpar}
\inferrule
  {\step{e}{e'}}
  {\step{\suc{e}}{\suc{e'}}}
\quad\textsc{(E-Succ)}

\inferrule
  { }
  {\step{\suc{i}}{i+1}}
\quad\textsc{(E-SuccInt)}

\inferrule
  {\step{e}{e'}}
  {\step{\prd{e}}{\prd{e'}}}
\quad\textsc{(E-Pred)}

\inferrule
  { }
  {\step{\prd{i}}{i-1}}
\quad\textsc{(E-PredInt)}

\inferrule
  {\step{e}{e'}}
  {\step{\iszero{e}}{\iszero{e'}}}
\quad\textsc{(E-IsZero)}

\inferrule
  { }
  {\step{\iszero{0}}{\true}}
\quad\textsc{(E-IsZeroZero)}

\inferrule
  {i \neq 0}
  {\step{\iszero{i}}{\false}}
\quad\textsc{(E-IsZeroNonZero)}
\end{mathpar}

\subsection*{String Operators}
\begin{mathpar}
\inferrule
  {\step{e_1}{e_1'}}
  {\step{\concat{e_1}{e_2}}{\concat{e_1'}{e_2}}}
\quad\textsc{(E-Concat1)}

\inferrule
  {\step{e_2}{e_2'}}
  {\step{\concat{s_1}{e_2}}{\concat{s_1}{e_2'}}}
\quad\textsc{(E-Concat2)}

\inferrule
  { }
  {\step{\concat{s_1}{s_2}}{s_1 \cdot s_2}}
\quad\textsc{(E-ConcatStr)}

\inferrule
  {\step{e}{e'}}
  {\step{\isemptystr{e}}{\isemptystr{e'}}}
\quad\textsc{(E-IsEmptyStr)}

\inferrule
  {|s| = 0}
  {\step{\isemptystr{s}}{\true}}
\quad\textsc{(E-IsEmptyStrTrue)}

\inferrule
  {|s| > 0}
  {\step{\isemptystr{s}}{\false}}
\quad\textsc{(E-IsEmptyStrFalse)}

\inferrule
  {\step{e}{e'}}
  {\step{\strlen{e}}{\strlen{e'}}}
\quad\textsc{(E-StrLen)}

\inferrule
  { }
  {\step{\strlen{s}}{|s|}}
\quad\textsc{(E-StrLenStr)}
\end{mathpar}

\section*{Typing Rules}

\begin{mathpar}
\inferrule
  { }
  {\typeof{\true}{\mbox{Bool}}}
\quad\textsc{(T-True)}

\inferrule
  { }
  {\typeof{\false}{\mbox{Bool}}}
\quad\textsc{(T-False)}

\inferrule
  { }
  {\typeof{i}{\mbox{Int}}}
\quad\textsc{(T-Int)}

\inferrule
  { }
  {\typeof{s}{\mbox{String}}}
\quad\textsc{(T-String)}

\inferrule
  {\typeof{e}{\mbox{Int}}}
  {\typeof{\suc{e}}{\mbox{Int}}}
\quad\textsc{(T-Succ)}

\inferrule
  {\typeof{e}{\mbox{Int}}}
  {\typeof{\prd{e}}{\mbox{Int}}}
\quad\textsc{(T-Pred)}

\inferrule
  {\typeof{e}{\mbox{Int}}}
  {\typeof{\iszero{e}}{\mbox{Bool}}}
\quad\textsc{(T-IsZero)}

\inferrule
  {\typeof{e_1}{\mbox{Bool}} \\
   \typeof{e_2}{T} \\
   \typeof{e_3}{T}}
  {\typeof{\ife{e_1}{e_2}{e_3}}{T}}
\quad\textsc{(T-If)}

\inferrule
  {\typeof{e_1}{\mbox{String}} \\
   \typeof{e_2}{\mbox{String}}}
  {\typeof{\concat{e_1}{e_2}}{\mbox{String}}}
\quad\textsc{(T-Concat)}

\inferrule
  {\typeof{e}{\mbox{String}}}
  {\typeof{\isemptystr{e}}{\mbox{Bool}}}
\quad\textsc{(T-IsEmptyStr)}

\inferrule
  {\typeof{e}{\mbox{String}}}
  {\typeof{\strlen{e}}{\mbox{Int}}}
\quad\textsc{(T-StrLen)}
\end{mathpar}

\section*{Why These Rules Give Progress and Preservation}

\textbf{Progress (simple statement).}
If $\typeof{e}{T}$, then either $e$ is already a value, or there exists $e'$ such that $\step{e}{e'}$.

\textbf{Why progress holds.}
The typing rules only allow valid inputs for each operator.
So a well-typed operator never gets stuck:
\begin{itemize}
  \item \texttt{succ}, \texttt{pred}, and \texttt{iszero} only receive Int expressions.
  \item \texttt{concat}, \texttt{isemptystr}, and \texttt{strlen} only receive String expressions.
  \item \texttt{if} always has a Bool condition.
\end{itemize}

\textbf{Preservation (simple statement).}
If $\typeof{e}{T}$ and $\step{e}{e'}$, then $\typeof{e'}{T}$.

\textbf{Why preservation holds.}
Each computation rule returns the correct result type:
\begin{itemize}
  \item \texttt{succ}/\texttt{pred}: Int $\rightarrow$ Int
  \item \texttt{iszero}: Int $\rightarrow$ Bool
  \item \texttt{concat}: String $\rightarrow$ String (with two String inputs)
  \item \texttt{isemptystr}: String $\rightarrow$ Bool
  \item \texttt{strlen}: String $\rightarrow$ Int
  \item \texttt{if}: result type is the common type of both branches
\end{itemize}
So one step of evaluation does not change the type.

\end{document}
